\scanpage{187}

Supposons maintenant la caractéristique nulle. Alors les représentations \inserted{lin.} rationnelles du \uline{groupe} \uline{abélien} $k$ correspondent à \uline{certaines} représentations lin. de l'\uline{algèbre de Lie abélienne} $k$ ; il résulte de la \deleted{\unsure{pr.}} prop. 1 qu'il s'agit exactement des représentations \uline{nilpotentes} de l'algèbre $k$, dont la classification est celle des $k[X]$-modules de type fini $X$-primaires. Ainsi :

\noindent\begin{tabular}{@{}|p{.92\linewidth}@{}}\uline{Prop. 2} La classification des repr. linéaires du groupe alg. abélien $k$ est, via les représentations infinitésimales associées, équivalente à la classification des $k[X]$-modules de type fini $(X)$-primaires.\end{tabular}

% slanted
\margin{$k$ de car $0$}

Les $k[X]$-modules $(X)$-primaires \uline{indécomposables} sont ceux qui sont du type $k[X]/X^{n}$ $(1\leq n<+\infty)$. On \deleted{\unsure{du}} voit aussitôt que le dual \inserted{sur $k$} du $k[X]$ module $E_n$ est isomorphe à $E_n$.

% slanted
\margin{La classification qu'on a donné des \uline{classes d'\unsure{opérateurs|invariants}} \uline{nilpotents} \deleted{\illegible{1}} dépend \uline{aussi} de la caractéristique.}

Ainsi

\uline{Corollaire} L'anneau \inserted{$\unsure{A|\mathcal A}(k)$} des classes de représentations du groupe alg. $k$ a pour base canonique

% archive sheet number: 186
